% Nouveau EGA V (cote 25): archive cote 25, site page 115 (PDF page 116).
% typescript. Draft compared: none.
% AI-transcribed by Claude Opus 5.5 (claude-opus-5-5, adaptive thinking, effort high), Message Batches API, batch ega25b (msgbatch_01A7n5YuXubUPNK1zsepi8RG), 2026-11-13; not yet checked by eye.
\scanpage{25/115}
\mspage{18}
changement de base en S; d'autre part, X-Z \inserted[typed]{(supposé ouvert, p ex T/S plat lpf.}\inserted{lorsque $T \to S$ est surjectif\unsure{)}} est universellement
schématiquement dense dans X rel à S. Ce dernier fait résulte du \xout{la compa-}
\xout{tibilité avec changement de base} par. 11 et du fait que pour tout $s \in S$,
$Z_s$ ne contient aucun point de $X_s$ associé à $\underline{O}_{X_s}$ (en effet, le support d'un
diviseur \deleted{de Cartier} sur $X_s$ ne contient aucun tel point).

Dans le cas où T \deleted{devient}\inserted{est} un fibré projectif $Q=P(F)$, $\mathrm{Div}_{Q/S}$
étant représentable par la somme des $P(\mathrm{Sym}^n(\check{F})) = P(n)$, on trouve
un morphisme $X-Z \longrightarrow P(n)$ . On dit\xout{\unsure{a}} que la famille de diviseurs
D est \uline{de degré} n, si le morphisme précédent se factorise par P(n); si
$X \neq \emptyset$ , d'où $X-Z \neq \emptyset$ , le n en question est bien déterminé par D. D'ail%
leurs pour définir cette notion, on n'a pas besoin à la rigueur du résul-%
tat de représentabilité énoncé\xout{s} plus haut, mais seulement d'avoir défini
les monomorphismes canoniques
\[ P(n) \longrightarrow \mathrm{Div}_{\check{P}/S} \]
(NB on pose $P = P(\check{F})$, d'où $Q = P(F) = \check{P}$ avec les notations des n°s précé-%
dents). \deleted{Cela pourra être fait à titre d'exemple dans le n° préconisé au
par. 19}. On appelle \uline{système linéaire de diviseurs sur} X/S , \uline{paramétré
par le fibré projectif} $Q = \check{P}$ , \uline{toute famille de diviseurs} \uline{sur} X/S
\uline{paramétrée par} Q, \uline{qui soit de degré} 1, i.e. \uline{définisse}
\[ \hand{f} : X-Z \longrightarrow P \qquad . \]
Ainsi, à un tel système \xout{\unsure{li}} linéaire de diviseurs\inserted{(et lorsque les fibres de $\check{P}$ sont tous $\neq \emptyset$)}, est associé canoniquement
une application rationelle de X dans un fibré projectif,\inserted{savoir un ``\deleted{\illegible{1}} \unsure{pseudo-morphisme} rel.: S''} \hand{(}en fait, un peu
mieux même qu'une application rationelle\hand{)}. Par construction même,
$D\,|(X-Z)x_S\check{P}$ n'est autre que l'image inverse par $f x id_{\check{P}}$ du diviseur
canonique (le diviseur d'in\unsure{cid}\xout{\unsure{e}}ence )\hand{\uline{H}} sur $Px_S\check{P}$ , Donc la connaissance
de $f: X-Z \longrightarrow P$ permet de réconstituer \xout{la fa} tout au moins la famil-%
le de diviseurs de X-Z induite par D , - donc si la famille est sans
\mspage{115}
